% ============================================================
% SECTION 2 — Related Work
% ============================================================
\section{Related Work}
\label{sec:related}

\paragraph{Spatial validation and leakage.}
Standard cross-validation assumes independent observations. When data
exhibit spatial autocorrelation, nearby locations leak information
across folds and inflate reported skill. Roberts et
al.~\cite{roberts2017cv} formalized spatial blocking strategies;
Ploton et al.~\cite{ploton2020spatial} demonstrated that large-scale
ecological mapping models can lose most of their apparent skill under
spatial validation; and Mil\`{a} et al.~\cite{mila2022nndm} proposed
nearest-neighbour distance matching for map-aware fold construction.
These methods address the prediction geometry. Our benchmark adopts
spatial blocking as the headline evaluation protocol and reports the
gap between random and blocked skill as the leakage diagnostic, but
extends the evaluation to five additional decision geometries that
spatial blocking alone does not address.

\paragraph{Scale, aggregation, and administrative-unit pitfalls.}
The Modifiable Areal Unit Problem~\cite{openshaw1984maup} established
that statistical relationships change with the spatial unit of
aggregation. Resolution mismatch between labels and predictions
introduces systematic error in remote sensing~\cite{workman2022resolution}.
In health and policy research, ZIP codes and Census-defined ZCTAs are
routinely conflated despite material spatial
mismatches~\cite{krieger2002zip, chen2025crosswalk}. Our substrate
operates at the ZCTA grain and documents the ZIP-as-ZCTA approximation
as a known limitation (C-001, \S\ref{sec:substrate}). Geographically
weighted regression~\cite{brunsdon1996gwr} provides a lens for
diagnosing spatial non-stationarity in model coefficients, which we
apply in the clustering and prediction geometries
(\S\ref{sec:spatial-diagnostics}).

\paragraph{Geospatial benchmarks.}
SustainBench~\cite{yeh2021sustainbench} evaluates sustainability
indicators across multiple tasks but reports per-task accuracy, not
per-decision sufficiency. SatCLIP~\cite{klemmer2023satclip} benchmarks
location embeddings; PDFM~\cite{agarwal2025pdfm} evaluates a
geospatial foundation model on population dynamics tasks with
state-holdout extrapolation. GeoYCSB~\cite{kim2023geoycsb} targets
database performance, not model evaluation.
GeoBenchX~\cite{krechetova2025geobenchx} and
GeoAnalystBench~\cite{zhang2025geoanalyst} benchmark agent-based
geospatial workflows. Dynabench~\cite{kiela2021dynabench} and
Ott et al.~\cite{ott2022saturation} identify benchmark saturation
as a systemic problem in ML evaluation. Each of these contributes a
dataset, a task, or a platform; none organizes evaluation around the
decision the prediction is used for or reports per-geometry evidence
sufficiency. That is the gap this benchmark addresses.

\paragraph{RSCT and GeoRSCT.}
The Representation--Solver Compatibility Theory
(RSCT)~\cite{martin2026rsct} decomposes model quality into a simplex
of signal ($R$), suppression ($S_\text{sup}$), and noise ($N$), with
derived measures for signal purity ($\alpha$), geometric compatibility
($\kappa$), and cross-fold stability ($\tau$, $\sigma$). GeoRSCT
(NeurIPS D\&B)~\cite{martin2026georsct} validates this machinery on a
27-task CONUS-scale geospatial benchmark. The present paper instantiates
the same certificate framework on a flood substrate organized around
decision geometries rather than regression targets, using the
certificate as an audit trail that explains verdict changes across
representation levels.
